\subsubsection{Complete seed-optimization calls}

The final measurements run serially after the tests, audit and publication
of the runtime change. Each experiment has five shuffled paired rounds,
old and new A/A controls, and one excluded warmup per arm and input. All
580 measured calls and 116 warmups complete. The source pins are identical
across the final audit and all three benchmarks.

The optimization experiment starts from a fresh stabilized-circle diagram,
constructs and independently checks its exterior, discovers its primitive
cocycle, runs the optimizer, checks the span witness, and validates the
resulting normal coordinates. The binary case first adds a global vertex
potential of amplitude $2^{4096}$, preserving the primitive class. The
baseline is the unchanged delivered solver with a native certificate adapter;
its own built-in check is retained, and the adapter also invokes the native
independent checker. The new producer already invokes that checker. The
common measurement endpoint additionally replays the returned span proof.
These are complete adapted seed-optimization calls, not isolated flow times.

\begin{center}
\small
\begin{tabular}{lrrrrr}
\toprule
Input & Old (ms) & New (ms) & Old/new & Old A/A & New A/A \\
\midrule
1 crossing & 8.582 & 7.315 & 1.187 & 1.032 & 0.987 \\
4 crossings & 72.640 & 51.040 & 1.437 & 1.000 & 1.000 \\
8 crossings & 238.285 & 152.483 & 1.557 & 0.987 & 0.998 \\
16 crossings & 895.695 & 499.241 & 1.787 & 1.005 & 0.985 \\
4 crossings, 4096 bits & 110.527 & 58.919 & 1.906 & 0.991 & 0.999 \\
\bottomrule
\end{tabular}

\end{center}

All one hundred measured calls and twenty warmups finish; driver elapsed
time is 25.414 seconds. Paired gains range from $1.187$ to $1.906$, beyond
the observed old A/A range $0.987$--$1.032$ and new A/A range
$0.985$--$1.000$. At sixteen crossings, the median falls from 895.695 ms
to 499.241 ms. The high-bit case falls from 110.527 ms to 58.919 ms.
Each implementation's certificate hash is stable across its repeats, and
the two optimum values agree. Different optimal potentials may produce
different valid witnesses.

\subsubsection{Full recognition with the native stage deliberately reached}

The second experiment disables earlier diagram, braid and invariant
shortcuts so that the source reaches the new native stage. Both arms use
the same current recognizer; only the span optimizer changes between the
delivered implementation and the native one. The raw-circle case exits
before either optimizer. The optimized-positive case returns a checked
normal-disc proof. The remaining cases continue through the exact
Khovanov fallback after the candidate queries fail. All verdicts agree with
the known inputs; no inconclusive recognition is treated as completed.

\begin{center}
\small
\begin{tabular}{lrrrrr}
\toprule
Input & Old (ms) & New (ms) & Old/new & Old A/A & New A/A \\
\midrule
Raw circle & 376.791 & 374.341 & 1.024 & 1.010 & 0.980 \\
Optimized positive & 299.693 & 288.459 & 1.034 & 0.982 & 0.959 \\
Genus-one miss & 84.344 & 85.157 & 0.986 & 1.003 & 1.002 \\
Trefoil & 92.423 & 85.683 & 1.064 & 1.057 & 0.998 \\
Figure-eight & 161.984 & 145.843 & 1.098 & 1.000 & 0.996 \\
\bottomrule
\end{tabular}

\end{center}

All one hundred measured recognitions and twenty warmups finish in
24.132 seconds of driver elapsed time. Paired ratios range from $0.986$
to $1.098$, much smaller than the optimization-stage gains. The
optimized-positive ratio $1.034$ is comparable to its new A/A fluctuation
($0.959$), and the trefoil's $1.064$ is close to its old A/A ratio $1.057$.
The raw-circle control follows the same code in both arms. These results
do not justify a broad whole-recognition speedup claim; source preparation,
component queries, proof replay and fallback can dominate the saved flow work.

\subsubsection{The ordinary compressed-group portfolio}

The last experiment compares release \code{535a914a6} with the new release
and the native option enabled, using the established compressed-group
portfolio on all nineteen ordinary cases. The earlier stages decide every
input, so there are zero native-seed attempts in all measured calls. This
is an integration and overhead control, not a measurement of native-seed
discovery or an expansion of automatic recognition coverage.

\begin{center}
\small
\begin{tabular}{lrrrrr}
\toprule
Input & Old (ms) & New (ms) & Old/new & Old A/A & New A/A \\
\midrule
survivor-00 & 15.021 & 15.090 & 0.995 & 0.994 & 0.999 \\
survivor-01 & 12.901 & 12.983 & 0.993 & 1.002 & 1.004 \\
survivor-02 & 14.435 & 14.373 & 1.035 & 1.020 & 1.000 \\
survivor-03 & 10.080 & 10.416 & 0.977 & 1.010 & 1.029 \\
mirror-03 & 22.651 & 22.667 & 1.000 & 0.928 & 0.990 \\
survivor-04 & 8.627 & 8.598 & 0.997 & 0.997 & 0.993 \\
survivor-05 & 17.494 & 17.499 & 1.002 & 1.005 & 0.997 \\
survivor-06 & 7.098 & 7.279 & 0.975 & 1.006 & 1.050 \\
survivor-07 & 18.311 & 18.548 & 0.993 & 1.003 & 1.001 \\
survivor-08 & 6.923 & 6.898 & 0.998 & 1.016 & 1.007 \\
mirror-08 & 18.854 & 19.028 & 0.989 & 1.000 & 1.005 \\
survivor-09 & 7.055 & 6.925 & 0.985 & 1.006 & 1.017 \\
survivor-10 & 6.436 & 6.354 & 1.021 & 1.020 & 0.988 \\
survivor-11 & 6.829 & 6.842 & 1.013 & 1.013 & 0.994 \\
gordian & 1223.718 & 1233.042 & 0.998 & 1.002 & 0.986 \\
Trefoil & 0.063 & 0.064 & 0.984 & 0.987 & 0.985 \\
figure\_eight & 0.068 & 0.070 & 0.972 & 0.988 & 1.016 \\
conway & 0.980 & 0.992 & 0.985 & 1.000 & 1.014 \\
kinoshita\_terasaka & 1.006 & 1.009 & 0.998 & 0.994 & 1.001 \\
\bottomrule
\end{tabular}

\end{center}

All 380 measured calls and 76 warmups finish; driver elapsed time is
34.309 seconds. Old/new paired ratios span $0.972$--$1.035$, old A/A ratios
$0.928$--$1.020$, and new A/A ratios $0.985$--$1.050$. Gordian's old and
new medians are 1.223718 and 1.233042 seconds, with paired ratio $0.998$.
No broad ordinary-workload gain or regression is established. The new
option remains disabled by default.
